% =========================================================================
\section{Planteamiento y estructura de la solución}\label{sec:estructura}
% =========================================================================

Todo el tema descansa sobre una sola propiedad: el operador diferencial es
\emph{lineal}. De ella se siguen tres hechos: que la solución general es
«homogénea más particular», que las soluciones de la homogénea forman un
espacio vectorial de dimensión $n$, y que una fuente complicada puede
descomponerse en impulsos elementales cuyas respuestas se suman
(\cref{sec:green}).

% -------------------------------------------------------------------------
\subsection{El operador lineal de orden \texorpdfstring{$n$}{n}}\label{sec:operador}
% -------------------------------------------------------------------------

\begin{definicion}[EDO lineal de orden $n$]\label{def:lineal}
Sean $I\subseteq\R$ un intervalo abierto y $a_0,\dots,a_n,b\in\mathcal C(I)$
con $a_n(x)\neq0$ para todo $x\in I$. Con $D\defeq\dif/\dif x$, el
\emph{operador diferencial lineal de orden $n$}
\begin{equation}\label{eq:operador}
  \Lop\defeq\sum_{j=0}^{n}a_j(x)\,D^{j}\colon\ \mathcal C^{n}(I)\longrightarrow\mathcal C(I)
\end{equation}
define la \emph{ecuación diferencial lineal de orden $n$}
\begin{equation}\label{eq:edo}
  \Lop[y]\equiv a_n(x)\,y^{(n)}+a_{n-1}(x)\,y^{(n-1)}+\dots+a_1(x)\,y'+a_0(x)\,y=b(x).
\end{equation}
Se llama \emph{homogénea} si $b\equiv0$ y \emph{completa} en caso
contrario; $b$ es la \emph{fuente} o término inhomogéneo.
\end{definicion}

\begin{proposicion}[Linealidad y superposición]\label{prop:linealidad}
Para cualesquiera $y_1,y_2\in\mathcal C^{n}(I)$ y $\alpha,\beta\in\C$,
\[
  \Lop[\alpha y_1+\beta y_2]=\alpha\,\Lop[y_1]+\beta\,\Lop[y_2].
\]
En consecuencia: \textup{(i)} el conjunto de soluciones de la homogénea es
el núcleo $\Ker\Lop$, un subespacio vectorial; \textup{(ii)} si
$\Lop[y_k]=b_k$ para $k=1,\dots,m$, entonces
$\Lop\bigl[\sum_k y_k\bigr]=\sum_k b_k$ (principio de superposición).
\end{proposicion}

\begin{proof}
La derivación es lineal y $a_j(x)$ multiplica punto a punto; (i) y (ii)
son consecuencias inmediatas.
\end{proof}

% -------------------------------------------------------------------------
\subsection{El espacio de soluciones tiene dimensión \texorpdfstring{$n$}{n}}\label{sec:dimension}
% -------------------------------------------------------------------------

\begin{teorema}[Existencia y unicidad; recordatorio]\label{teo:eyu}
Para cada $x_0\in I$ y cada $(\eta_0,\dots,\eta_{n-1})\in\R^{n}$ existe una
única $y\in\mathcal C^{n}(I)$ solución de \eqref{eq:edo} con
$y^{(j)}(x_0)=\eta_j$, $j=0,\dots,n-1$
\textup{\cite[cap.~3]{CODDINGTON}, \cite[cap.~IV]{TEMARIO}}.
\end{teorema}

\begin{teorema}[Dimensión del núcleo]\label{teo:dimension}
$\dim\Ker\Lop=n$. Más precisamente, para cada $x_0\in I$ la aplicación
«datos iniciales»
\begin{equation}\label{eq:Phi}
  \Phi_{x_0}\colon\Ker\Lop\longrightarrow\R^{n},\qquad
  \Phi_{x_0}(y)=\bigl(y(x_0),\,y'(x_0),\,\dots,\,y^{(n-1)}(x_0)\bigr),
\end{equation}
es un isomorfismo de espacios vectoriales.
\end{teorema}

\begin{proof}
$\Phi_{x_0}$ es lineal porque la derivación y la evaluación lo son. Aplicando
el \cref{teo:eyu} con $b\equiv0$: la \emph{existencia} dice que todo vector
de datos es imagen de alguna solución, luego $\Phi_{x_0}$ es sobreyectiva;
la \emph{unicidad} dice que esa solución es única, luego $\Phi_{x_0}$ es
inyectiva. Así $\Ker\Lop\cong\R^{n}$.
\end{proof}

\begin{observacion}[¿$n$ soluciones o $n$ constantes?]\label{obs:n-constantes}
Ambas afirmaciones expresan lo mismo. Conviene distinguir tres objetos: el
conjunto de soluciones $\Ker\Lop$, que es infinito; una \emph{base}
$\{\hat y_1,\dots,\hat y_n\}$ o \emph{sistema fundamental}, las piezas
fijas; y las $n$ constantes $c_1,\dots,c_n$, que son las coordenadas de una
solución concreta en esa base, $y=c_1\hat y_1+\dots+c_n\hat y_n$. Es la misma
situación que en $\R^{3}$: infinitos vectores, tres vectores de la base y
tres coordenadas para cada uno. Fijar una solución equivale a dar $n$
números, sean los datos iniciales o las constantes $c_i$
(\cref{fig:isomorfismo}).
\end{observacion}

\begin{ejemplo}\label{ej:oscilador-base}
Para $y''+y=0$ ($n=2$), $\{\cos x,\sin x\}$ es una base y la solución
general $y=c_1\cos x+c_2\sin x$ tiene dos constantes. En $x_0=0$,
$y(0)=c_1$ e $y'(0)=c_2$: aquí las constantes son directamente los datos
iniciales, porque $\Phi_0(\cos)=(1,0)$ y $\Phi_0(\sin)=(0,1)$ son la base
canónica de $\R^{2}$.
\end{ejemplo}

\begin{notacion}[Subespacio generado]\label{not:generado}
El \emph{subespacio generado} por $f_1,\dots,f_k$ es el conjunto de sus
combinaciones lineales, que denotamos indistintamente
\begin{equation}\label{eq:generado}
  \gen{f_1,\dots,f_k}=\spanop\{f_1,\dots,f_k\}
  =\Bigl\{\,\sum_{i=1}^{k}c_if_i\ \Big|\ c_i\in\C\Bigr\}.
\end{equation}
Por ejemplo, fijado $n\in\N$, $\{1,x,x^{2},\dots,x^{n}\}$ es un sistema de
generadores del subespacio $\R_n[x]\subset\R[x]$ de polinomios de
coeficientes reales y grado $\le n$:
\[
  \gen{1,x,x^{2},\dots,x^{n}}=\R_n[x]=\Bigl\{\,\sum_{i=0}^{n}a_ix^{i}\ \Big|\ a_i\in\R\Bigr\}
\]
(y $\C_n[x]$ si los coeficientes son complejos). Así, si
$\{\hat y_1,\dots,\hat y_n\}$ es un sistema fundamental,
$\Ker\Lop=\gen{\hat y_1,\dots,\hat y_n}=\spanop\{\hat y_1,\dots,\hat y_n\}$.
\end{notacion}

% Figura: el isomorfismo «datos iniciales» y una familia de soluciones
\begin{figure}[ht]
\centering
\begin{tikzpicture}
  % ---------------- (a) Esquema del isomorfismo ----------------
  \begin{scope}
    \node[font=\small, text=texto, anchor=west] at (-0.3,3.05) {(a) $\Phi_{x_0}\colon\Ker\Lop\to\R^{n}$};
    % espacio de soluciones: caja con curvas
    \draw[cian, line width=0.7pt, rounded corners=6pt, fill=cajafondo] (0,0) rectangle (2.7,2.5);
    \node[font=\small, text=cian, anchor=north] at (1.35,-0.12) {$\Ker\Lop$};
    \draw[magenta, line width=0.9pt] plot[domain=0.25:2.45, samples=60, smooth] (\x,{1.25+0.55*sin(deg(2.6*\x))});
    \draw[ambar, line width=0.9pt] plot[domain=0.25:2.45, samples=60, smooth] (\x,{1.25+0.45*cos(deg(2.6*\x))});
    \draw[verde, line width=0.9pt] plot[domain=0.25:2.45, samples=60, smooth] (\x,{1.25+0.6*sin(deg(2.6*\x+40))*exp(-0.3*\x)});
    % R^n: ejes
    \begin{scope}[shift={(5.4,0.35)}]
      \draw[eje] (0,0) -- (2.2,0);
      \draw[eje] (0,0) -- (0,2.1);
      \draw[eje] (0,0) -- (-0.75,-0.55);
      \draw[-{Stealth[length=6pt,width=5pt]}, line width=1pt, ambar] (0,0) -- (1.35,1.25);
      \fill[ambar] (1.35,1.25) circle (1.3pt);
      \node[font=\small, text=cian, anchor=north] at (0.75,-0.47) {$\R^{n}$};
      \node[etiqgris, anchor=west, align=left] at (1.42,1.42) {$(y(x_0),\dots,$\\$\ y^{(n-1)}(x_0))$};
    \end{scope}
    % flechas
    \draw[flecha, line width=0.8pt, cian] (2.85,1.7) to[bend left=18]
      node[above, font=\scriptsize, text=gris] {$\Phi_{x_0}$} (5.15,1.7);
    \draw[flecha, line width=0.8pt, violeta] (5.15,0.95) to[bend left=18]
      node[below, font=\scriptsize, text=gris, align=center] {$\Phi_{x_0}^{-1}$} (2.85,0.95);
    \node[font=\scriptsize, text=gris, align=center] at (4.0,-0.55) {existencia $\Rightarrow$ sobre\\unicidad $\Rightarrow$ inyectiva};
  \end{scope}
  % ---------------- (b) Familia de soluciones de y''+y=0 ----------------
  \begin{axis}[at={(9.9cm,0)}, anchor=south west, oscuro, width=5.75cm, height=4.4cm,
      xmin=-0.6, xmax=6.6, ymin=-1.75, ymax=1.75, xtick={0,3.1416,6.2832},
      xticklabels={$0$,$\pi$,$2\pi$}, ytick={-1,0,1}, xlabel={$x$},
      title={(b) $y''+y=0$: datos en $x_0=0$}, samples=120, domain=-0.6:6.6]
    \addplot[c1] {cos(deg(x))};
    \addplot[c2] {sin(deg(x))};
    \addplot[c3] {cos(deg(x))+sin(deg(x))};
    \addplot[c4] {-0.5*cos(deg(x))+0.8*sin(deg(x))};
    \addplot[only marks, mark=*, mark size=1.6pt, blanco] coordinates {(0,1) (0,0) (0,-0.5)};
  \end{axis}
\end{tikzpicture}
\caption{(a) La aplicación $\Phi_{x_0}$ asigna a cada solución de la
homogénea sus $n$ datos iniciales en $x_0$; el teorema de existencia y
unicidad la convierte en un isomorfismo, de modo que $\dim\Ker\Lop=n$.
(b) Soluciones de $y''+y=0$ con datos $(y(0),y'(0))=(1,0)$, $(0,1)$,
$(1,1)$ y $(-\tfrac12,\tfrac45)$: cada par de datos determina una única curva.}
\label{fig:isomorfismo}
\end{figure}


% -------------------------------------------------------------------------
\subsection{Independencia lineal y wronskiano}\label{sec:wronskiano}
% -------------------------------------------------------------------------
No basta con tener $n$ soluciones: deben ser independientes. Por ejemplo,
$\{\cos x,\,2\cos x\}$ sólo genera una recta de $\Ker\Lop$.

\begin{definicion}[Matriz wronskiana]\label{def:wronskiano}
Dadas $\hat y_1,\dots,\hat y_n\in\mathcal C^{n-1}(I)$, su \emph{matriz
wronskiana} y su \emph{wronskiano} son
\begin{equation}\label{eq:wronskiana}
  W(x)\defeq
  \begin{pmatrix}
    \hat y_1 & \cdots & \hat y_n\\
    \hat y_1' & \cdots & \hat y_n'\\
    \vdots & & \vdots\\
    \hat y_1^{(n-1)} & \cdots & \hat y_n^{(n-1)}
  \end{pmatrix}\!(x),
  \qquad w(x)\defeq\det W(x).
\end{equation}
\end{definicion}

\begin{proposicion}[Criterio del wronskiano]\label{prop:wronskiano}
Sean $\hat y_1,\dots,\hat y_n\in\Ker\Lop$. Son equivalentes:
\begin{enumerate}[label=\textup{(\alph*)}]
  \item $\{\hat y_1,\dots,\hat y_n\}$ es una base de $\Ker\Lop$;
  \item $w(x_0)\neq0$ para \emph{algún} $x_0\in I$;
  \item $w(x)\neq0$ para \emph{todo} $x\in I$.
\end{enumerate}
\end{proposicion}

\begin{proof}
\emph{Lectura matricial de $W(x_0)$.} La columna $i$ de $W(x_0)$ es
precisamente $\Phi_{x_0}(\hat y_i)$. Multiplicar $W(x_0)$ por un vector
$\bm c=(c_1,\dots,c_n)^{\T}\in\R^{n}$ es combinar sus columnas, y por la
linealidad de $\Phi_{x_0}$,
\begin{equation}\label{eq:Wc}
  W(x_0)\,\bm c
  =\begin{pmatrix}
     \sum_i c_i\,\hat y_i(x_0)\\ \vdots\\ \sum_i c_i\,\hat y_i^{(n-1)}(x_0)
   \end{pmatrix}
  =\sum_{i=1}^{n}c_i\,\Phi_{x_0}(\hat y_i)
  =\Phi_{x_0}(y),\qquad y\defeq\sum_{i=1}^{n}c_i\,\hat y_i\in\Ker\Lop .
\end{equation}
Así, resolver $W(x_0)\,\bm c=\bm 0$ es buscar combinaciones de las $\hat y_i$
con datos iniciales nulos en $x_0$.

\emph{(a)$\Rightarrow$(c).} Supongamos $w(x_0)=0$ para algún $x_0\in I$.
Entonces existe $\bm c\neq\bm 0$ con $W(x_0)\,\bm c=\bm 0$ y, por
\eqref{eq:Wc}, $y=\sum_ic_i\hat y_i$ cumple
$y(x_0)=y'(x_0)=\dots=y^{(n-1)}(x_0)=0$. La función nula también resuelve
$\Lop[y]=0$ con esos datos en $x_0$; por la \emph{unicidad} del
\cref{teo:eyu} (inyectividad de $\Phi_{x_0}$), $y\equiv0$ en $I$. Es una
combinación lineal nula con $\bm c\neq\bm 0$, lo que contradice que
$\{\hat y_i\}$ sea base.

\emph{(c)$\Rightarrow$(b).} Inmediato.

\emph{(b)$\Rightarrow$(a).} Si $\sum_ic_i\hat y_i\equiv0$ en $I$, también se
anulan en $I$ sus derivadas hasta orden $n-1$, en particular en $x_0$; por
\eqref{eq:Wc}, $W(x_0)\,\bm c=\bm 0$ y, como $W(x_0)$ es invertible,
$\bm c=\bm 0$. Las $\hat y_i$ son, pues, $n$ vectores independientes de
$\Ker\Lop$, que tiene dimensión $n$ (\cref{teo:dimension}): forman una base.
\end{proof}

\begin{idea}[Papel de la unicidad]
Antes de usar la unicidad, en (a)$\Rightarrow$(c) sólo se sabe algo
\emph{en el punto} $x_0$: $y$ y sus $n-1$ primeras derivadas se anulan
allí. Eso no basta, porque $x^{3}$ se anula en $0$ junto con sus dos primeras derivadas sin
ser la función nula, y la independencia lineal es una afirmación sobre
\emph{todo} $I$. La unicidad del problema de valores iniciales es la que
propaga el anulamiento de $x_0$ a todo el intervalo, y sólo está disponible
porque las $\hat y_i$ resuelven una misma ecuación lineal con $a_n\neq0$.
La implicación (b)$\Rightarrow$(a), en cambio, vale para funciones
cualesquiera de $\mathcal C^{n-1}(I)$.
\end{idea}

\begin{observacion}[El recíproco falla sin ecuación]\label{obs:wronskiano-general}
Para funciones \emph{arbitrarias}, $w\not\equiv0$ implica independencia
lineal (el argumento de (b)$\Rightarrow$(a) no usa la ecuación), pero no al
revés. Sean $f(x)=x^{2}$ y $g(x)=x\abs{x}$ en $\R$.
\begin{enumerate}[label=(\roman*)]
  \item \emph{Son independientes.} Si $\alpha f+\beta g\equiv0$, evaluando
        en $x=1$ y $x=-1$ se obtiene $\alpha+\beta=0$ y $\alpha-\beta=0$, luego
        $\alpha=\beta=0$.
  \item \emph{Su wronskiano es nulo.} $g\in\mathcal C^{1}(\R)$ con
        $g'(x)=2\abs{x}$, así que
        $w(x)=x^{2}\cdot2\abs{x}-x\abs{x}\cdot2x=0$ para todo $x\in\R$.
  \item \emph{Dónde falla la demostración.} Con $\bm c=(1,-1)^{\T}$,
        \[
          y=f-g=\begin{cases}0, & x\ge0,\\ 2x^{2}, & x<0,\end{cases}
        \]
        tiene $y(0)=y'(0)=0$, pero $y\not\equiv0$. No hay unicidad que lo
        impida: $f$ y $g$ no resuelven ninguna ecuación
        $y''+p(x)\,y'+q(x)\,y=0$ con $p,q\in\mathcal C(\R)$, pues entonces la
        \cref{prop:wronskiano} exigiría $w\neq0$ en todo $\R$.
\end{enumerate}
\end{observacion}

\begin{ejemplo}[Dos wronskianos \cite{TEMARIO}]\label{ej:wronskianos}
\leavevmode
\begin{enumerate}[label=(\roman*)]
  \item $x^{3}y'''-3x^{2}y''+6x\,y'-6y=0$ en $I=(0,\infty)$ admite las
        soluciones $x$, $x^{2}$, $x^{3}$, y
        \[
          w(x)=\begin{vmatrix}x & x^{2} & x^{3}\\ 1 & 2x & 3x^{2}\\ 0 & 2 & 6x\end{vmatrix}=2x^{3}\neq0 ,
        \]
        luego la solución general es $y=k_1x+k_2x^{2}+k_3x^{3}$.
  \item Para $y''+y=0$, las soluciones $\cos x$, $\sin x$, $\eul^{\iu x}$ son
        dependientes, $\eul^{\iu x}-\cos x-\iu\sin x=0$ (como debe ser: tres
        vectores en un espacio de dimensión $2$), y en efecto su wronskiano
        $3\times3$ es nulo.
\end{enumerate}
\end{ejemplo}

La dicotomía de la \cref{prop:wronskiano} («$w$ no se anula nunca o es
idénticamente nulo») admite una versión cuantitativa: el wronskiano de $n$
soluciones satisface una EDO lineal de primer orden. Necesitamos derivar un
determinante.

\begin{lema}[Derivada de un determinante]\label{lem:derivada-det}
Sea $A(x)=\bigl(a_{kl}(x)\bigr)$ una matriz $n\times n$ de funciones
derivables en $I$. Entonces $\det A$ es derivable y
\begin{equation}\label{eq:derivada-det}
  (\det A)'=\sum_{k=1}^{n}\det A_k ,
\end{equation}
donde $A_k$ se obtiene sustituyendo la fila $k$ de $A$ por su derivada.
\end{lema}

\begin{proof}
Por la fórmula de Leibniz,
$\det A=\sum_{\sigma\in S_n}\operatorname{sgn}\sigma\,\prod_{k=1}^{n}a_{k\sigma(k)}$,
una suma finita de productos de funciones derivables. Por la regla del
producto,
\[
  (\det A)'=\sum_{\sigma\in S_n}\operatorname{sgn}\sigma\sum_{k=1}^{n}
  a_{1\sigma(1)}\cdots a_{k\sigma(k)}'\cdots a_{n\sigma(n)}
  =\sum_{k=1}^{n}\ubr{\sum_{\sigma\in S_n}\operatorname{sgn}\sigma\,
  a_{1\sigma(1)}\cdots a_{k\sigma(k)}'\cdots a_{n\sigma(n)}}{\det A_k},
\]
donde se ha intercambiado el orden de los dos sumatorios finitos.
\end{proof}

\begin{teorema}[Fórmula de Abel--Liouville]\label{teo:abel}
Sean $\hat y_1,\dots,\hat y_n\in\Ker\Lop$ y $x_0\in I$. Para todo $x\in I$,
\begin{equation}\label{eq:abel}
  \boxed{\;w(x)=w(x_0)\,\exp\!\Bigl(-\int_{x_0}^{x}\frac{a_{n-1}(s)}{a_n(s)}\,\dif s\Bigr)\;}
\end{equation}
En particular, $w$ o es idénticamente nulo en $I$ o no se anula en ningún
punto de $I$.
\end{teorema}

\begin{proof}
Numeramos las filas de $W$ por el orden de derivación $j=0,\dots,n-1$; la
fila $j$ es $\bigl(\hat y_1^{(j)},\dots,\hat y_n^{(j)}\bigr)$. Como
$\hat y_i\in\mathcal C^{n}(I)$, todas las entradas de $W$ son de clase
$\mathcal C^{1}$, y por el \cref{lem:derivada-det}
$w'=\sum_{j=0}^{n-1}\det W_j$, donde $W_j$ tiene la fila $j$ sustituida por
$\bigl(\hat y_1^{(j+1)},\dots,\hat y_n^{(j+1)}\bigr)$.
\begin{enumerate}[label=\textcolor{cian}{\arabic*.}, leftmargin=1.5em, itemsep=3pt]
  \item \emph{Filas $j\le n-2$.} La fila nueva coincide con la fila $j+1$
        de $W$; $W_j$ tiene dos filas iguales y $\det W_j=0$.
  \item \emph{Fila $j=n-1$.} Sobrevive un único término:
        \[
          w'=\det W_{n-1}=
          \begin{vmatrix}
            \hat y_1 & \cdots & \hat y_n\\
            \vdots & & \vdots\\
            \hat y_1^{(n-2)} & \cdots & \hat y_n^{(n-2)}\\
            \hat y_1^{(n)} & \cdots & \hat y_n^{(n)}
          \end{vmatrix}.
        \]
  \item \emph{Uso de la ecuación.} Como $\Lop[\hat y_i]=0$ y $a_n\neq0$,
        $\hat y_i^{(n)}=-\sum_{l=0}^{n-1}\frac{a_l}{a_n}\,\hat y_i^{(l)}$: la
        última fila de $W_{n-1}$ es la combinación
        $-\sum_{l}\frac{a_l}{a_n}\cdot(\text{fila }l\text{ de }W)$. Por la
        linealidad del determinante en esa fila, $\det W_{n-1}$ es
        $-\sum_l\frac{a_l}{a_n}$ por el determinante de $W$ con la última fila
        sustituida por la fila $l$. Para $l\le n-2$ hay dos filas iguales y
        el sumando es nulo; para $l=n-1$ se recupera $W$. Así
        \begin{equation}\label{eq:abel-edo}
          w'(x)=-\frac{a_{n-1}(x)}{a_n(x)}\,w(x).
        \end{equation}
  \item \emph{Integración.} Sea
        $P(x)\defeq\int_{x_0}^{x}a_{n-1}(s)/a_n(s)\,\dif s$, de clase
        $\mathcal C^{1}$ con $P'=a_{n-1}/a_n$. Por \eqref{eq:abel-edo},
        $\bigl(w\,\eul^{P}\bigr)'=\bigl(w'+P'w\bigr)\eul^{P}=0$, luego
        $w\,\eul^{P}$ es constante en el intervalo $I$ e igual a su valor en
        $x_0$, que es $w(x_0)$ porque $P(x_0)=0$. Esto es \eqref{eq:abel}.
        La última afirmación se sigue de que la exponencial no se anula.\qedhere
\end{enumerate}
\end{proof}

\begin{observacion}\label{obs:abel}
\leavevmode
\begin{enumerate}[label=(\roman*)]
  \item El \cref{teo:abel} da otra prueba de (b)$\Rightarrow$(c) en la
        \cref{prop:wronskiano}. La equivalencia con (a) sigue necesitando la
        unicidad.
  \item \emph{Comprobación} en el \cref{ej:wronskianos}\,(i):
        $a_{2}/a_{3}=-3x^{2}/x^{3}=-3/x$, luego
        $w(x)=w(x_0)\exp\int_{x_0}^{x}\frac{3}{s}\,\dif s=w(x_0)\,(x/x_0)^{3}$,
        que con $w(x_0)=2x_0^{3}$ da $w(x)=2x^{3}$. Para $y''+y=0$
        ($a_1=0$), el wronskiano es constante:
        $w=\cos^{2}x+\sin^{2}x=1$.
  \item Con coeficientes constantes,
        $w(x)=w(x_0)\,\eul^{-(a_{n-1}/a_n)(x-x_0)}$.
\end{enumerate}
\end{observacion}

% -------------------------------------------------------------------------
\subsection{Estructura de la solución general}\label{sec:general}
% -------------------------------------------------------------------------

\begin{teorema}[Homogénea más particular]\label{teo:estructura}
Sean $\{\hat y_1,\dots,\hat y_n\}$ un sistema fundamental e $y_p$ una
solución cualquiera de \eqref{eq:edo}. Entonces toda solución de
\eqref{eq:edo} es de la forma
\begin{equation}\label{eq:general}
  y(x)=\ubr{\sum_{i=1}^{n}c_i\,\hat y_i(x)}{y_h\,\in\,\Ker\Lop}+\;y_p(x),
  \qquad c_1,\dots,c_n\in\R .
\end{equation}
\end{teorema}

\begin{proof}
Por linealidad, $\Lop[y-y_p]=b-b=0$, luego $y-y_p\in\Ker\Lop$, que es el
espacio $\gen{\hat y_1,\dots,\hat y_n}$ (\cref{not:generado}).
\end{proof}

El trabajo se reparte así en dos problemas (\cref{fig:mapa}): hallar
$y_h$, lo que haremos para coeficientes constantes en las
\cref{sec:constantes,sec:multiples}, y hallar $y_p$, lo que haremos en
general en las \cref{sec:variacion,sec:green}.

% Figura: mapa del tema (y = y_h + y_p)
\begin{figure}[ht]
\centering
\begin{tikzpicture}
  \node[nodocian, align=center, text width=3.6cm] (eq) at (0,0)
    {$\Lop[y]=b$\\[1pt]\footnotesize\color{gris} lineal, orden $n$};
  \node[nodoambar, align=left, text width=5.3cm] (yh) at (6.2,1.05)
    {\textcolor{ambar}{\bfseries $y_h\in\Ker\Lop$}\ \ \footnotesize\color{gris}($\dim=n$)\\[2pt]
     \footnotesize coef.\ constantes: $p(\lambda)=0$ (\S\ref{sec:constantes}),\\
     \footnotesize raíces múltiples (\S\ref{sec:multiples})};
  \node[nodovioleta, align=left, text width=5.3cm] (yp) at (6.2,-1.05)
    {\textcolor{violeta}{\bfseries $y_p$ particular}\\[2pt]
     \footnotesize variación de constantes (\S\ref{sec:variacion})\\
     \footnotesize $\Longleftrightarrow$ función de Green (\S\ref{sec:green})};
  \node[nodo, draw=blanco, align=center, text width=2.9cm] (y) at (11.6,0)
    {$y=y_h+y_p$\\[1pt]\footnotesize\color{gris} $n$ constantes libres};
  \draw[flecha, cian] (eq.east) -- ++(0.45,0) |- (yh.west);
  \draw[flecha, cian] (eq.east) -- ++(0.45,0) |- (yp.west);
  \draw[flecha, ambar] (yh.east) -| ($(y.north)+(-0.3,0)$);
  \draw[flecha, violeta] (yp.east) -| ($(y.south)+(-0.3,0)$);
\end{tikzpicture}
\caption{Estructura del tema: por la linealidad, el problema se separa en
hallar el núcleo de $\Lop$ y una solución particular.}
\label{fig:mapa}
\end{figure}

